\documentclass[a4paper,12pt]{jsarticle}
\usepackage[margin=25mm]{geometry}
\usepackage{amsmath,amssymb}
\usepackage{enumitem}

\begin{document}

\begin{center}
{\LARGE 数I・A 共通テスト本番レベル（やや難）解答}
\end{center}

\vspace{1em}

\section*{第1問}
\[
f(x)=x^2-2ax+a+2=(x-a)^2-a^2+a+2
\]
\begin{enumerate}[label=\arabic*.]
  \item 最小値は
  \[
  -a^2+a+2.
  \]
  \item $f(x)\le0$ の解集合の長さは、判別式を $D$ とすると $\sqrt{D}$。
  \[
  D=(-2a)^2-4(a+2)=4(a^2-a-2).
  \]
  長さ $\ge2$ より
  \[
  \sqrt{D}\ge2 \iff D\ge4
  \iff a^2-a-3\ge0.
  \]
  よって
  \[
  a\le\frac{1-\sqrt{13}}{2}\ \text{または}\ a\ge\frac{1+\sqrt{13}}{2}.
  \]
  \item $0<\alpha<1<\beta$ は $f(1)<0$ かつ2実根をもつことと同値。
  \[
  f(1)=1-2a+a+2=3-a<0 \iff a>3.
  \]
  また $a>3$ なら判別式は正で、
  \[
  \alpha\beta=a+2>0,\quad \alpha+\beta=2a>0
  \]
  より2根は正。
  よって
  \[
  a>3.
  \]
  \item $a=4$ のとき
  \[
  f(x)=x^2-8x+6=(x-4)^2-10.
  \]
  
  \[
  |f(x)|\le3
  \iff -3\le (x-4)^2-10\le3
  \iff 7\le (x-4)^2\le13.
  \]
  よって
  \[
  x\in[4-\sqrt{13},\,4-\sqrt7]\cup[4+\sqrt7,\,4+\sqrt{13}].
  \]
\end{enumerate}

\section*{第2問}
\begin{enumerate}[label=\arabic*.]
  \item
  \[
  \binom{12}{4}=495.
  \]
  \item 偶数6枚、奇数6枚より
  \[
  \frac{\binom62\binom62}{\binom{12}4}
  =\frac{225}{495}=\frac{5}{11}.
  \]
  \item 条件「最大値が10」では、10を含み残り3枚を1～9から選ぶ。
  \[
  \text{条件事象}=\binom93=84.
  \]
  さらに最小値が3なら、3と10を含み、残り2枚を4～9から選ぶので
  \[
  \text{有利事象}=\binom62=15.
  \]
  よって条件付き確率は
  \[
  \frac{15}{84}=\frac{5}{28}.
  \]
  \item 1～12を3で割った余り0,1,2の個数はそれぞれ4個。
  和が3の倍数となる組数を数えると
  \[
  1+96+16+16+36=165
  \]
  （内訳：$(4,0,0),(2,1,1),(1,3,0),(1,0,3),(0,2,2)$）。
  よって
  \[
  \frac{165}{495}=\frac13.
  \]
\end{enumerate}

\section*{第3問}
\begin{enumerate}[label=\arabic*.]
  \item 余弦定理より
  \[
  BC^2=6^2+10^2-2\cdot6\cdot10\cos120^\circ
  =36+100+60=196.
  \]
  よって
  \[
  BC=14.
  \]
  \item 面積
  \[
  S=\frac12\cdot6\cdot10\sin120^\circ
  =30\cdot\frac{\sqrt3}{2}=15\sqrt3.
  \]
  \item 半周長を $s$ とすると
  \[
  s=\frac{6+10+14}{2}=15.
  \]
  内接円半径 $r$ は $S=rs$ より
  \[
  r=\frac{S}{s}=\frac{15\sqrt3}{15}=\sqrt3.
  \]
  \item 中線の長さ公式
  \[
  AM=\frac12\sqrt{2AB^2+2AC^2-BC^2}
  =\frac12\sqrt{2\cdot6^2+2\cdot10^2-14^2}
  =\frac12\sqrt{76}=\sqrt{19}.
  \]
\end{enumerate}

\section*{第4問}
\begin{enumerate}[label=\arabic*.]
  \item データは20個。第10番目と第11番目はいずれも55なので
  \[
  \text{中央値}=55.
  \]
  \item 下位10個の中央値は第5・6番目（50,50）の平均より
  \[
  Q_1=50.
  \]
  上位10個の中央値は第15・16番目（65,65）の平均より
  \[
  Q_3=65.
  \]
  よって
  \[
  IQR=Q_3-Q_1=15.
  \]
  \item 対称性または計算より
  \[
  \bar{x}=\frac{40\cdot1+45\cdot2+50\cdot3+55\cdot4+60\cdot4+65\cdot3+70\cdot2+75\cdot1}{20}
  =\frac{1150}{20}=57.5.
  \]
  \item
  \[
  \sum (x-\bar{x})^2
  =306.25+312.5+168.75+25+25+168.75+312.5+306.25
  =1625.
  \]
  よって
  \[
  s^2=\frac{1625}{20}=81.25=\frac{325}{4},
  \quad
  s=\sqrt{\frac{325}{4}}=\frac{5\sqrt{13}}{2}.
  \]
\end{enumerate}

\end{document}
